% ==========================================================================
%  Chapter 108 -- Theoretical Compendium of the Branch-Selection Programme
% ==========================================================================

\chapter{Theoretical Compendium of the Branch-Selection Programme}
\label{ch:theoretical-compendium}

\paragraph{Normative status.}
Reference chapter. It fixes, with explicit definitions and the exact
mathematics implemented, every concept the branch-selection and
multiscale programme of Caps.~\ref{ch:branch-selection-reversals}--%
\ref{ch:kinematic-regularization} relies on, so that no later reading
of the campaigns or of the code admits ambiguity. Everything stated
here is implemented in the repository; where fall\_n deliberately
extends a published formulation, the extension is stated as such.

% ─────────────────────────────────────────────────────────────────────────
\section{Notation and conventions}
\label{sec:tc-notation}

\begin{definition}[Tensor storage]
Second-order symmetric tensors are stored in Voigt order
$\{11,22,33,23,13,12\}$ as $\mathbb{R}^6$ vectors; engineering shear
strains carry the factor 2. Where orthogonal transformations of the
constitutive matrix are needed (crack systems), the implementation
converts to the Mandel basis (off-diagonal entries scaled by
$\sqrt2$), applies the rotation, and converts back, so that
$\mathbf{C}^{\text{mandel}} = \mathbf{Q}\,\mathbf{C}\,\mathbf{Q}^\top$
preserves the energy inner product.
\end{definition}

\begin{definition}[Stress invariants]
\label{def:tc-invariants}
With principal stresses $\sigma_1\le\sigma_2\le\sigma_3$, deviators
$s_i = \sigma_i - \sigma_o$ and $J_2 = \tfrac12\sum_i s_i^2$,
\[
  \sigma_o = \tfrac13\,\sigma_{kk}, \qquad
  \tau_o^{\text{oct}} = \sqrt{\tfrac{2J_2}{3}}, \qquad
  \cos 3\theta = \tfrac{3\sqrt3}{2}\, J_3\, J_2^{-3/2},
\]
where $\tau_o^{\text{oct}}$ is the octahedral shear stress of
\cite{KoBathe2026} (Eq.~2b there). The historical fall\_n
transcription computed $\tau_o^{\text{code}} = \sqrt{2J_2}/3 =
\tau_o^{\text{oct}}/\sqrt3$; the branch
\texttt{fix/kobathe-tau-octahedral} restores the octahedral
definition. Every statement in this compendium holds for either
definition provided it is used consistently; the quantitative
constants of Section~\ref{sec:tc-kobathe} are calibrated for the
octahedral one (Cap.~\ref{ch:kinematic-regularization},
fidelity audit).
\end{definition}

\begin{definition}[Pseudo-time and protocol]
The quasi-static protocol is parameterised by pseudo-time
$p\in(0,1]$, discretised into $N_T$ targets $p_k = k/N_T$. For the
reference column, $N_T = N_{\text{pre}} + N_{\text{lat}}$ with
$N_{\text{pre}}=4$ axial-preload targets and
$N_{\text{lat}} = 2\,S\,n_{\text{seg}}$ lateral targets ($S$ steps
per segment, substep factor 2, $n_{\text{seg}}=12$ segments from four
amplitudes $\pm50/100/150/200$\,mm, each $0\to+A\to-A\to0$). The
prescribed tip drift is the piecewise-linear map
$\delta(p)$ of the cyclic protocol; a \emph{reversal} is a sign
change of $\delta(p_{k+1})-\delta(p_k)$.
\end{definition}

\begin{definition}[Base shear convention]
$V$ denotes the horizontal support resultant. The CSVs store the
\emph{reaction} (negative under positive push); figures plot
$V=-\text{reaction}$ so that loops run bottom-left to top-right.
\end{definition}

% ─────────────────────────────────────────────────────────────────────────
\section{The incremental problem and the frozen-state potential}
\label{sec:tc-incremental}

\begin{definition}[Committed state, trial state, ratchets]
The material history $\mathcal{H}$ (plastic strains, crack systems,
loading maxima) advances only at \emph{commit}, executed once per
accepted step through
\texttt{accept\_external\_solution\_step}/\texttt{commit\_trial\_state}.
Within a step every residual or tangent evaluation is a \emph{trial}:
it reads $\mathcal{H}$ but does not write it. Two \emph{ratchets}
(monotone history variables evaluated over committed configurations
only) drive the model of Section~\ref{sec:tc-kobathe}: the loading
maximum $\tau_o^{\max}$ and the per-crack strain maximum
$\max_\tau \bar e^{C}_n$.
\end{definition}

\begin{proposition}[Frozen-state conservativity]
\label{prop:tc-conservative}
Within a step (fixed $\mathcal{H}$, crack count frozen by
\texttt{check\_new\_cracks=false}), the residual field is the gradient
of an incremental potential, $\mathbf R(\mathbf u)=\nabla\Pi(\mathbf
u)$, because every constitutive branch then in force is a secant
(path-independent) law. Consequently the line integral
\[
  \Pi(\mathbf u)-\Pi(\mathbf u_{\mathrm{ref}})
  \;=\; \int_0^1 \mathbf R\!\left(\mathbf u_{\mathrm{ref}} +
        s\,\Delta\mathbf u\right)^{\!\top} \Delta\mathbf u \; ds,
  \qquad \Delta\mathbf u=\mathbf u-\mathbf u_{\mathrm{ref}},
\]
is path-independent and is evaluated exactly (for polynomial
integrands) by the three-point Gauss--Legendre rule of
\texttt{NLAnalysisEnergyBackend}. Energies are comparable only within
one step and from one $\mathbf u_{\mathrm{ref}}$.
\end{proposition}

\begin{remark}
The hypothesis matters: across a commit the ratchets advance and
$\Pi$ changes branch, so no global potential exists over the protocol.
All energy-based criteria below are strictly per-step.
\end{remark}

% ─────────────────────────────────────────────────────────────────────────
\section{Regularized Newton continuation and its selection layers}
\label{sec:tc-lm}

\begin{definition}[LM step]
At control target $p$, warm-started at the last accepted $\mathbf u$,
each iteration solves
\[
  (\mathbf K + \mu\,\mathbf I)\, d\mathbf u = -\mathbf R(\mathbf u),
\]
with $\mathbf K$ the assembled tangent and $\mu\ge0$ the
Levenberg--Marquardt regularizer \cite{Levenberg1944,Marquardt1963}.
The schedule multiplies $\mu$ by \texttt{grow} on rejection and by
\texttt{drop} on acceptance, capped at
$\mu_{\max}=f_{\mu}\lVert\operatorname{diag}\mathbf K\rVert$; this is
the classical trust-region correspondence (larger $\mu$
$\Leftrightarrow$ smaller trust radius)
\cite{NocedalWright2006,ConnGouldToint2000}. Acceptance requires a
monotone residual decrease; a stagnation cutoff accepts the best
regularized iterate, and a step is reported converged if
$\lVert\mathbf R\rVert$ meets the tolerances or the acceptance floor.
The secant predictor extrapolates the warm start along $p$ with a
bounded scale and falls back to the plain warm start if it raises the
residual.
\end{definition}

\begin{proposition}[Descent property]
\label{prop:tc-descent}
If $\mathbf K+\mu\mathbf I$ is positive definite then
$\mathbf R^\top d\mathbf u = -\mathbf R^\top(\mathbf K+\mu\mathbf
I)^{-1}\mathbf R < 0$; with
Proposition~\ref{prop:tc-conservative}, $d\mathbf u$ is a descent
direction of $\Pi$. In the large-$\mu$ limit $d\mathbf u \to
-\mu^{-1}\nabla\Pi$, i.e.\ gradient descent, which cannot converge to
a maximiser or saddle of $\Pi$.
\end{proposition}

\begin{definition}[Energy line search]
\label{def:tc-ls}
Gate \texttt{energy\_linesearch}: a candidate step $\alpha\,d\mathbf
u$ (starting at $\alpha=1$) is accepted only if the residual test of
the LM step holds \emph{and}
$\Pi(\mathbf u+\alpha\,d\mathbf u)\le\Pi(\mathbf u)$; otherwise
$\alpha$ is halved up to a fixed number of backtracks, and if no
$\alpha$ achieves energy non-increase the iteration is \emph{rejected}
so that $\mu$ grows. By Proposition~\ref{prop:tc-descent} the
rejection loop terminates in the gradient-descent regime, which
excludes convergence to energy maxima (the mechanism the double-well
unit test isolates). The rule prevents basin jumps; it does not claim
global descent across residual ridges.
\end{definition}

\begin{definition}[Deflation]
For known roots $\{\mathbf u_i\}$, the deflation operator
\cite{FarrellBirkissonFunke2015}
$\eta(\mathbf u)=\prod_i\left(\lVert\mathbf u-\mathbf
u_i\rVert^{-p}+\alpha\right)$ modifies the Newton direction by the
rank-one Sherman--Morrison factor
$\tau = \left(1-\frac{\nabla\eta\cdot d\mathbf
u}{\eta}\right)^{-1}$ (clamped), and the acceptance test compares
deflated residuals $\lVert\mathbf R\rVert\,\eta$. The campaign verdict
(Cap.~\ref{ch:driver-integration}) is negative for this material:
repulsion selects \emph{some other} root, with no control over its
quality; the implementation retains it as a gated experiment with a
never-worse-than-base fallback.
\end{definition}

\begin{definition}[Energy trust-region step]
The hybrid TAO step minimises the incremental potential with the
bound-constrained Newton trust region \texttt{bntr} of PETSc/TAO
\cite{BalayPETScManual}: objective $\Pi$, gradient $\mathbf R$,
Hessian $\mathbf K$ (Proposition~\ref{prop:tc-conservative}); the LM
iterate seeds it, the accepted minimiser is handed back. It selects by
construction (descent to a minimiser) and is fragile exactly where
$\Pi$ loses $C^1$ regularity; the two-physics matrix of
Cap.~\ref{ch:kinematic-regularization} bounds its domain of use.
\end{definition}

% ─────────────────────────────────────────────────────────────────────────
\section{The constitutive model, as implemented}
\label{sec:tc-kobathe}

The concrete model follows \cite{KoBathe2026}: strain-driven
plastic-fracturing with (i) loading/unloading classification by
$\tau_o \gtrless \tau_o^{\max}$; (ii) empirically softened bulk/shear
moduli $K_s(\sigma_o), G_s(\tau_o)$ and their tangential companions
$K_c, G_c$ (Eq.~10 and Appendix~A there, coefficients
$A,b,c,d,k,l,m,n$ as functions of $f_c$); (iii) an internal confining
pressure $\sigma_{id}$; (iv) up to three orthogonal smeared cracks per
point with orientation from the largest principal stress, opening
ratchet $0<\bar e^{C}_n\le\max_\tau \bar e^{C}_n$, branch matrices
$\mathbf C^{C(n)}$, and crushing (all-zero law) beyond three; (v)
crack-band regularization through $G_F$ and a length scale $l_b$.
Unloading is incrementally linear in the softened moduli.

fall\_n's deliberate extensions, each motivated by the reversal
pathologies the source never exercises
(Cap.~\ref{ch:kinematic-regularization} fidelity audit): softened
moduli evaluated at the historical maxima (irreversible,
damage-secant unloading without hysteresis), a $\tanh$-smoothed crack
closure replacing the discrete branch switch in the
\texttt{production\_stabilized} profile, opening-graduated shear
retention and Poisson decoupling on the crack planes, floors
$0.01\,K_e, 0.01\,G_e$ on the softened moduli, and the delay-damage
cap on per-step opening growth.

% ─────────────────────────────────────────────────────────────────────────
\section{Element-wise FE\textsuperscript{2} enrichment}
\label{sec:tc-fe2}

\begin{definition}[Coupling site and RVE $=$ macro element]
\label{def:tc-site}
A \emph{coupling site} attaches to one macro beam element $e$ (a
Gauss point of it carries the injected section law). Its local model
is a full three-dimensional discretisation of the \emph{entire domain
of that element}: the prismatic solid between the element's end
sections (\texttt{align\_to\_beam(endpoint\_A, endpoint\_B)}), with
the complete cross-section, embedded reinforcement, and persistent
material history. Nothing smaller than the element is ever
sub-modelled. In the reference column (eight $0.4$\,m elements) one
site enriches the base element; the container
(\texttt{MultiscaleModel}) holds arbitrarily many (site, local) pairs
--- the building drivers enrich many elements at once --- so the
scheme extends to frames by registering one RVE per structural
element wherever local physics matters, with feedback active where
warranted. Enriching \emph{all} elements of the column is a
configuration, not a code change.
\end{definition}

\begin{definition}[Downscaling: kinematic transfer]
Let $\hat{\mathbf u}(\xi), \hat{\bm\theta}(\xi)$ be the macro
element's interpolated section translation and rotation. The transfer
imposes on boundary nodes of the RVE the rigid-section map
\[
  \mathbf u(y,z)\big|_{\xi}
  \;=\; \hat{\mathbf u}(\xi) + \hat{\bm\theta}(\xi)\times(0,y,z)^\top ,
\]
as Dirichlet data. Mode \texttt{faces} applies it at the two end
sections ($\xi=\pm1$) only, leaving the interior free; mode
\texttt{layers} (B1) applies it additionally at the element-layer
stations $\xi_k=-1+2k/n_z$, a stronger control that removes local
basin ambiguity at a measured cost in global fidelity
(over-constraint, Cap.~\ref{ch:kinematic-regularization}).
\end{definition}

\begin{definition}[Upscaling: section operator]
The local model returns the pair $(\mathbf f, \mathbf D)$ of
section resultants and tangent. The production operator computes
$\mathbf f\in\mathbb{R}^6$ as the reduction of the boundary reactions
of the loaded face to $[N,M_y,M_z,V_y,V_z,M_t]$, and
$\mathbf D\in\mathbb{R}^{6\times6}$ by static condensation of the
local stiffness onto the six generalized section coordinates. The
volume-average operator (comparison path) replaces $\mathbf f$ by
resultants of the Gauss-point mean stress; used alone it yields an
inconsistent pair (flagged
\texttt{forces\_consistent\_with\_tangent=false}) whose live feedback
is unstable (Cap.~\ref{ch:energy-linesearch-fe2}); a consistent
volume operator requires its own tangent (finite differences of the
averaged forces) and remains future work.
\end{definition}

\begin{definition}[Affine injection]
The macro consumes the feedback as the frozen affine section law
$\mathbf f(\bm\varepsilon) = \mathbf f_0 + \mathbf
D\,(\bm\varepsilon-\bm\varepsilon_{\mathrm{ref}})$ at the site's
Gauss point, where $\bm\varepsilon_{\mathrm{ref}}$ is the macro
section strain at which $(\mathbf f_0,\mathbf D)$ was homogenised. A
stale $\bm\varepsilon_{\mathrm{ref}}$ extrapolates; the frozen-history
spiral of Cap.~\ref{ch:energy-linesearch-fe2} was precisely an
ancient anchor extrapolating across cycles.
\end{definition}

\begin{definition}[Staggered fixed point, relaxation, metrics]
One coupled step iterates: macro trial solve under the current
injection $\to$ per site, restore the local start-of-step state,
impose the trial kinematics, solve, homogenise $\to$ relax the
response against the previous iterate $\to$ test convergence. The
force and tangent metrics are relative gaps
\[
  \mathrm{res}_F = \frac{\lVert \mathbf f^{(k)} - \mathbf
  f^{(k-1)}\rVert}{\max(\lVert\mathbf f^{(k)}\rVert, s_F)},
  \qquad
  \mathrm{res}_K = \frac{\lVert \mathbf D^{(k)} - \mathbf
  D^{(k-1)}\rVert_F}{\max(\lVert\mathbf D^{(k)}\rVert_F, s_K)} ,
\]
against the staggered tolerance (3\% here). Relaxation policies
implement
$\mathbf x^{(k)} \leftarrow \omega\,\mathbf x^{(k)} +
(1-\omega)\,\mathbf x^{(k-1)}$ blended consistently on
$(\mathbf f,\mathbf D)$ at a common $\bm\varepsilon_{\mathrm{ref}}$:
constant $\omega$, or Aitken $\Delta^2$ dynamic relaxation
\cite{KuttlerWall2008}
\[
  \omega_{k+1} = -\,\omega_k\,
  \frac{\Delta\mathbf r_k \cdot (\Delta\mathbf r_{k+1}-\Delta\mathbf
  r_k)}{\lVert \Delta\mathbf r_{k+1}-\Delta\mathbf r_k \rVert^2},
\]
clamped to $[\omega_{\min},\omega_{\max}]$ and reset at each step's
first iteration (single-site state; per-site instances required
before frame-scale use).
\end{definition}

\begin{definition}[Recovery regimes and engagement]
A step that meets the staggered tolerance commits in
\emph{strict} regime. Otherwise the hybrid observation window
re-solves the macro under the last converged (or cleared) feedback,
commits, then \emph{evolves and commits the locals} under the final
macro state and refreshes the feedback anchor
(\texttt{evolve\_locals\_in\_hybrid}; the staggered-failure path
falls into this same recovery). Engagement policy: coupling may start
at any step; the measured hierarchy
(Cap.~\ref{ch:kinematic-regularization}) is \emph{virgin engagement
at the first post-yield reversal} $>$ short one-way warmup $>$
immediate $>$ long one-way warmup, because the unloading branch is
elastic and single-valued (consistency at engagement is free) whereas
one-way-tracked history is not equilibrium-consistent with the
coupled system.
\end{definition}

The scheme instantiates the classical FE\textsuperscript{2} programme
\cite{Feyel1999,FeyelandChaboche2000,Feyel2003,Miehe2002,Kouznetsova2004,Geers2010}
in its element-wise, structural-section variant: the ``unit cell'' is
not a periodic microstructure but the element's own domain, the
macro-to-micro link is section kinematics rather than deformation
gradients, and the upscaled quantity is a section law rather than a
stress. All well-posedness levers identified in that literature
(relaxation, operator consistency, engagement) reappear here in
section form.

% ─────────────────────────────────────────────────────────────────────────
\section{The metaheuristic layer}
\label{sec:tc-ca}

The Cultural Algorithm \cite{Reynolds1994Cultural,ReynoldsPeng2004}
maximises objectives $f:\mathbb{R}^d\to\mathbb{R}$ over a box, with a
population evolved under belief-space influence and elitism; a single
seeded \texttt{mt19937\_64} makes runs bit-reproducible. The
implementation keeps the framework's META character explicit
\cite{Maheri2021CulturalSurvey}: candidate solutions carry
\emph{traits} (not a ``genome'' --- that vocabulary belongs to one
instantiation), the population space is an interchangeable
\texttt{PopulationModel} policy (evolutionary programming by default;
a BLX-$\alpha$ genetic-recombination model is provided as the
alternative reading), and the belief space aggregates any subset of
the literature's five knowledge sources --- normative (elite
intervals), situational (best exemplar), domain (fitness-weighted
elite consensus), history (temporal trend of generation bests) and
topographic (per-trait occupancy cells). On the multimodal Rastrigin
benchmark the five-source belief space improves the two-source result
($-0.87$ vs $-3.01$ at equal budget). In the solver-tuning instance
the traits are the eight-dimensional LM configuration (log-scales for
$\mu_0$ and the acceptance floor; rounded integers for stagnation and
subdivision), the objective replays a checkpointed protocol window,
and the fitness is
$w_1 n_{\text{conv}}/m - w_2\max(0, |V|_{\text{peak}}/V_{\text{ref}}
- 1) - w_3\bar\imath - w_4\,\mathrm{zigzag}$ with $w=(1,1,0.2,w_4)$,
where $V_{\text{ref}}$ is the run's own pre-reversal envelope in the
v1 fitness and the \emph{physical plateau of the structural element}
(Cap.~\ref{ch:cultural-program}) in the v2 fitness. The same layer is
reserved for control-device optimisation
\cite{LaraValenciaEchavarriaValencia2024}.

% ─────────────────────────────────────────────────────────────────────────
\section{Diagnostic metrics}
\label{sec:tc-metrics}

\begin{definition}[Campaign metrics]
For a protocol trace $\{(d_k, V_k, c_k)\}_{k=1}^{N}$ (drift, base
shear, converged flag):
\emph{noconv} $=\#\{k: c_k=0\}$ (floor- or stagnation-accepted
steps); peak $|V|_{\max}=\max_k |V_k|$; envelope variables
$V_{\mathrm{exc}}$ (running max of $|V|$ over converged steps of the
current excursion), $V_{\mathrm{env}}$ (its value frozen at each
reversal; reference of the spike detector
$|V|> s\,V_{\mathrm{env}}$), $V_{\mathrm{all}}$ (historical max);
\emph{roughness} $\Sigma\Delta V = \sum_k |V_{k}-V_{k-1}|$ (total
variation; compare per-step when $N$ differs); \emph{zigzag} $=$ the
fraction of protocol-monotone steps ($ (d_k-d_{k-1})(d_{k-1}-d_{k-2})
> 0$) at which $\Delta V$ alternates sign
($(V_k-V_{k-1})(V_{k-1}-V_{k-2})<0$) --- the sawtooth measure, with
the uncoupled fiber column at $2.2\%$ as the smooth reference.
\end{definition}

\begin{remark}
Boundedness ($|V|_{\max}$, spike counts) and smoothness (zigzag,
roughness) are independent qualities; the anti-sawtooth campaign of
Cap.~\ref{ch:kinematic-regularization} shows limiter-based
boundedness coexisting with fixed-point sawtooth, and cures for one
degrading the other (the 2\% force clamp).
\end{remark}
